\documentclass[12pt]{article}
\usepackage{fullpage}
\usepackage{amsmath,amssymb,amsthm}

\newcommand{\cL}{\mathcal L}
\newcommand{\Dr}{\operatorname{Dr}}
\newcommand{\Trop}{\operatorname{Trop}}
\newcommand{\rk}{\operatorname{rk}}
\newcommand{\cl}{\operatorname{cl}}
\newcommand{\op}{\operatorname{op}}

\newtheorem{theorem}{Theorem}
\newtheorem{lemma}{Lemma}
\newtheorem{proposition}{Proposition}
\newtheorem{criterion}{Criterion}

\begin{document}

\begin{center}
{\bf A polar transform for the relative Dressian}
\end{center}

\section*{0. Conventions and local inputs}

Let $E$ be the ground set, let $r=\rk M$, and let $\mu_M$ be the
trivial rank $r$ valuated matroid supported on the bases of $M$.  I use
the standard model of the relative Dressian in which a rank $k$ point is
a projective valuated matroid $\nu$ with $\nu\preceq\mu_M$, where
$\preceq$ is the valuated quotient relation.  By
Brandt--Eur--Zhang, this is the same as inclusion of projective tropical
linear spaces:
\[
        \Trop(\theta)\subseteq \Trop(\nu)
        \quad\Longleftrightarrow\quad \theta\preceq\nu .
\]
For valuated matroids $p,q$ of ranks $a\le b$, the quotient relation is
equivalent to the full incidence-Plucker relations
\[
 \min_{i\in J\setminus I}\{p(I\cup i)+q(J\setminus i)\}
 \quad\hbox{attained at least twice}                         \tag{A}
\]
for all $I\in\binom E{a-1}$ and $J\in\binom E{b+1}$, together with the
Plucker relations for $p$ and $q$.  All-$\infty$ minima are interpreted
as attained at every displayed index.

The construction below is completely formal except for the following
local-to-global point.  It is known in the adjacent-rank case and in the
ordinary case, and it is the precise remaining place where a proof or
reference is still needed for arbitrary rank gaps.

\begin{criterion}[two-rank local quotient criterion]\label{crit:local}
Let $p,q$ be valuated matroids of ranks $a<b$ whose ordinary supports
are in quotient order.  If the nested incidence relations
\[
 \min_{i\in T\setminus S}\{p(S\cup i)+q(T\setminus i)\}
 \quad\hbox{attained at least twice}                         \tag{N}
\]
hold for all $S\subseteq T$, $|S|=a-1$, $|T|=b+1$, then the full
relations {\rm (A)} hold.
\end{criterion}

I also use the following standard one-point completion.  If
$\nu\preceq\mu_M$ has rank $k$, define
\[
 h_j(B)=
 \begin{cases}
 \min\{\nu(C): B\subseteq C,\ |C|=k\}, & j\le k,\ |B|=j,\\
 \min\{\nu(C): C\subseteq B,\ |C|=k\}, & j\ge k,\ B\in\mathcal I_j(M),\\
 \infty, & j\ge k,\ B\notin\mathcal I_j(M).
 \end{cases}                                                \tag{1}
\]
Murota's truncation and relative elongation theorems, equivalently the
Dress--Terhalle well-layering construction, give valuated matroids
$h_j$ and adjacent quotients $h_j\preceq h_{j+1}$; the top layer is
projectively $\mu_M$.  For the endpoint, if $B'=B-e+f$ are adjacent
$M$-bases and $C\subseteq B$ minimizes the second line of (1), then
either $e\notin C$ and $C\subseteq B'$, or the incidence relation for
$\nu\preceq\mu_M$ with lower set $C-e$ and upper set $B\cup f$ gives a
$k$-subset of $B'$ of value at most $\nu(C)$.  Hence the value of
$h_r$ is constant on the basis graph.

Finally, in the proof that the polar of a single layer is again a
valuated matroid, I use the complete-flag local criterion: a complete
ordinary flag support together with all adjacent three-term tropical
incidence relations is a complete valuated flag.  This is the
non-uniform, $\infty$-valued version of the local exchange theorem for
valuated flag matroids.  A fully self-contained proof of this criterion
is listed as an open item at the end.

\section*{1. Consequences of the flat-lattice duality}

\begin{lemma}\label{lem:modular}
If $L=\cL(M)$ admits an order-reversing involution
$x\mapsto x^\perp$, then $L$ is modular and
$\rk(x^\perp)=r-\rk(x)$.
\end{lemma}

\begin{proof}
The anti-automorphism sends $\hat0$ to $\hat1$ and reverses maximal
chains, so it sends rank $i$ flats to rank $r-i$ flats.  The flat
lattice is semimodular.  Applying the anti-automorphism to the
semimodular inequality gives the reverse inequality; hence equality
holds for every pair of flats.
\end{proof}

Thus $(x\vee y)^\perp=x^\perp\wedge y^\perp$ and
$(x\wedge y)^\perp=x^\perp\vee y^\perp$.

\begin{lemma}[flat constancy]\label{lem:flatconst}
Let $\nu\preceq \mu_M$ have rank $k$.  If $B,B'$ are independent
$k$-subsets of $M$ with $\cl(B)=\cl(B')$, then $\nu(B)=\nu(B')$.
Also $\nu(A)=\infty$ whenever $A$ is dependent in $M$.
\end{lemma}

\begin{proof}
The support matroid of $\nu$ is an ordinary quotient of $M$, so its
bases are independent in $M$.  The case $k=0$ is trivial.  For $k>0$,
the basis graph of the restriction of $M$ to a rank $k$ flat is
connected, so it is enough to compare $B=A\cup b$ and $B'=A\cup b'$
with $|A|=k-1$ and $\cl(Ab)=\cl(Ab')$.  Choose a common complement
$C$ so that $AbC$ is an $M$-basis; since $Ab$ and $Ab'$ span the same
rank $k$ flat, $Ab'C$ is also an $M$-basis.  In the incidence relation
for $\nu\preceq\mu_M$ with lower set $A$ and upper set
$A\cup\{b,b'\}\cup C$, the only finite $\mu_M$-terms are the $b$ and
$b'$ terms, because deleting an element of $C$ leaves the dependent set
$Abb'$.  Hence $\nu(Ab)=\nu(Ab')$.
\end{proof}

Write $\bar\nu(X)$ for the common value on independent $k$-sets
spanning the rank $k$ flat $X$, and put $\bar\nu(X)=\infty$ when no
basis of $\nu$ spans $X$.

\section*{2. The polar transform and supports}

For a rank $k$ quotient $\nu\preceq\mu_M$ define a rank $r-k$ vector by
\[
 \nu^{\perp_M}(A)=
 \begin{cases}
   \bar\nu\bigl((\cl A)^\perp\bigr),
      & A\text{ independent in }M,\ |A|=r-k,\\
   \infty,&\text{otherwise.}
 \end{cases}                                                 \tag{2}
\]
This is compatible with projective rescaling, and applying (2) twice
recovers the original flat coordinates.

\begin{lemma}[ordinary polar rank]\label{lem:ordinary}
Let $N$ be an ordinary quotient of $M$ of rank $k$, with rank function
$\rho_N$.  The function
\[
        \rho^\perp(A)=\rk_M(A)-k+\rho_N((\cl A)^\perp)          \tag{3}
\]
is a matroid rank function.  Its bases of size $r-k$ are precisely the
finite support of (2).  If $N_p$ is a quotient of $N_q$, then
$N_q^{\perp_M}$ is a quotient of $N_p^{\perp_M}$.
\end{lemma}

\begin{proof}
First take $A=X$ a flat.  Since $N$ is a quotient of $M$, rank
increments in $N$ are bounded by those in $M$.  Hence
$\rho^\perp(X)\ge0$ follows from
$k-\rho_N(X^\perp)\le r-\rk(X^\perp)=\rk(X)$, and
$\rho^\perp(X)\le\rk(X)$ is immediate from $\rho_N(X^\perp)\le k$.
If $X\le Y$, then
\[
\rho^\perp(Y)-\rho^\perp(X)
 =\rk(Y)-\rk(X)-\bigl(\rho_N(X^\perp)-\rho_N(Y^\perp)\bigr)
\]
is nonnegative and is at most $\rk(Y)-\rk(X)$.  For a cover it is
therefore $0$ or $1$.  Submodularity on flats follows from modularity of
$\rk_M$ and submodularity of $\rho_N$, because $^\perp$ interchanges
meet and join.  Extending by $A\mapsto\cl(A)$ gives a submodular,
monotone, unit-increment rank function on all subsets, and
$\rho^\perp(A)\le\rk_M(A)\le |A|$.

If $A$ is $M$-independent, then
$\rho^\perp(A)=|A|-k+\rho_N((\cl A)^\perp)$, so an $M$-independent set
of size $r-k$ is a basis exactly when
$\rho_N((\cl A)^\perp)=k$.  Dependent sets of $M$ have
$\rho^\perp(A)<|A|$.

Now suppose $N_p$ has rank $a$, $N_q$ has rank $b$, and
$N_p$ is a quotient of $N_q$.  For flats $X\le Y$ the above formula
gives
\[
\rho_{N_q^{\perp_M}}(Y)-\rho_{N_q^{\perp_M}}(X)
 \le
\rho_{N_p^{\perp_M}}(Y)-\rho_{N_p^{\perp_M}}(X),
\]
because the $N_p$-increment on $Y^\perp\le X^\perp$ is at most the
$N_q$-increment.  This is the strong-map rank-increment criterion; for
arbitrary subsets one replaces them by their $M$-closures.
\end{proof}

\begin{lemma}[elementary opposition]\label{lem:elementary}
Let $p\preceq q\preceq\mu_M$, with $\rk p=d$ and $\rk q=d+1$.  Then
$q^{\perp_M}$ and $p^{\perp_M}$ satisfy all adjacent local three-term
relations of ranks $r-d-1$ and $r-d$.
\end{lemma}

\begin{proof}
Fix $A$ of size $r-d-2$ and distinct $x,y,z\notin A$.  Put
$U=\cl(A)$, $W=\cl(Axyz)$, and $\delta=\rk W-\rk U$.  If $A$ is
dependent, or if $\delta\le1$, all three transformed terms are
$\infty$.

If $\delta=2$, let $G=W^\perp$ and $H=U^\perp$, of ranks $d$ and
$d+2$.  A finite term has common $p$-coordinate $\bar p(G)$ and
$q$-coordinate one of
$\bar q((\cl(Ax))^\perp)$, $\bar q((\cl(Ay))^\perp)$,
$\bar q((\cl(Az))^\perp)$.  In the rank-two interval $[U,W]$, loops
give only $\infty$ terms and parallel points give equal finite
coordinates.  In the remaining case the three points are distinct.
Choose a basis $S$ of $G$, representatives of the three dual points in
$[G,H]$, and a common complement extending $H$ to an $M$-basis.  The
incidence relation for $q\preceq\mu_M$ has exactly these three finite
$q$-terms; adding $\bar p(G)$ gives the desired relation.

If $\delta=3$, then $x,y,z$ form a rank-three frame over $U$.  Let
$G=W^\perp$ and
\[
 P_x=(\cl(Ayz))^\perp,\quad P_y=(\cl(Axz))^\perp,\quad
 P_z=(\cl(Axy))^\perp .
\]
These are the three atoms over $G$ in the dual frame, and
$(\cl(Ax))^\perp=P_y\vee P_z$, with the analogous identities for
$y,z$.  Choose a basis $S$ of $G$ and representatives
$a_e\in P_e\setminus G$.  The adjacent incidence relation for
$p\preceq q$ applied to $S$ and $a_x,a_y,a_z$ is exactly the transformed
three-term relation.
\end{proof}

\begin{proposition}\label{prop:set}
Assuming the complete-flag local criterion from Section 0, the formula
$\nu\mapsto\nu^{\perp_M}$ defines an involution of the set underlying
$\Dr(M)$.
\end{proposition}

\begin{proof}
Insert $\nu$ into the complete flag (1).  Lemma \ref{lem:ordinary}
reverses the ordinary supports, and Lemma \ref{lem:elementary} gives
all adjacent three-term relations for
$h_r^{\perp_M},h_{r-1}^{\perp_M},\ldots,h_0^{\perp_M}$.  The endpoints
are the trivial rank $0$ valuation and $\mu_M$.  The complete-flag
criterion makes this reversed sequence a complete valuated flag, so in
particular $\nu^{\perp_M}\preceq\mu_M$.  Involutivity is immediate
from (2).
\end{proof}

\section*{3. Conditional order reversal}

\begin{lemma}[circuit minima]\label{lem:circuit}
Let $w\preceq\mu_M$ have rank $k\ge1$, let $A$ be an independent
$(k-1)$-set of $M$, and let $C$ be a circuit of $M/A$ with $|C|\ge2$.
Then $\min_{i\in C} w(A\cup i)$ is attained at least twice.
\end{lemma}

\begin{proof}
Let $F=\cl(A\cup C)$.  For every $i\in C$, the set $A\cup(C-i)$ is
independent and spans $F$.  Choose one complement $D$ of $F$ to an
$M$-basis.  In the incidence relation for $w\preceq\mu_M$ with lower
set $A$ and upper set $A\cup C\cup D$, the finite $\mu_M$-terms are
exactly the terms indexed by $i\in C$.
\end{proof}

\begin{theorem}\label{thm:order}
Assume Proposition \ref{prop:set} and Criterion \ref{crit:local}.  If
$p\preceq q\preceq\mu_M$, then
\[
        q^{\perp_M}\preceq p^{\perp_M}.
\]
\end{theorem}

\begin{proof}
Let $\rk p=a$ and $\rk q=b$.  Proposition \ref{prop:set} gives
valuated matroids, and Lemma \ref{lem:ordinary} gives the ordinary
quotient order.  It remains, by Criterion \ref{crit:local}, to check
nested incidence for $q^{\perp_M}\preceq p^{\perp_M}$.

Fix $A\subseteq T$ with
$|A|=r-b-1$, $|T|=r-a+1$, and put $B=T\setminus A$,
$m=|B|=b-a+2$.  If $A$ is dependent, all terms are $\infty$.  Otherwise
let $U=\cl(A)$, $W=\cl(T)$.  A term can be finite only if $T-i$ is
independent; hence if $\delta=\rk W-\rk U\le m-2$, all terms are
$\infty$.

Suppose $\delta=m$.  Then $T$ is independent.  Let $G=W^\perp$ and
$H=U^\perp$, of ranks $a-1$ and $b+1$.  For $i\in B$ set
$X_i=(\cl(T-i))^\perp$ and $Y_i=(\cl(Ai))^\perp$.  The flats
$\cl(A\cup J)$, $J\subseteq B$, form a Boolean frame in $[U,W]$; after
applying $^\perp$, the $X_i$ are the atoms over $G$ and
$Y_i=\bigvee_{j\ne i}X_j$.  Choosing a basis $S$ of $G$ and
representatives $x_i\in X_i\setminus G$, the nested relation for
$p\preceq q$ with lower set $S$ and upper set $S\cup\{x_i:i\in B\}$
has exactly the transformed terms.

It remains to consider $\delta=m-1$.  In $M/A$, the set $B$ has nullity
one and hence a unique circuit $C$.  If $C$ is a loop, all terms are
$\infty$.  For $i\notin C$, the set $T-i$ is dependent.  For $i\in C$,
the set $T-i$ is independent and spans $W$, so the
$p^{\perp_M}(T-i)$ factor is the common value $\bar p(W^\perp)$.
Lemma \ref{lem:circuit}, applied to $w=q^{\perp_M}$ and the circuit
$C$ of $M/A$, shows that the remaining minimum is attained at least
twice.
\end{proof}

\section*{4. Agreement with flats}

Let $F$ be a flat of rank $f$, and let $q_F$ be the trivial rank
$r-f$ valuated matroid of $M/F\oplus U_{0,F}$.  For an independent
$(r-f)$-set $B$ with $Z=\cl(B)$,
\[
        q_F(B)<\infty\quad\Longleftrightarrow\quad Z\vee F=\hat1.
\]
For an independent $f$-set $A$ with $Y=\cl(A)$, formula (2) gives
\[
 q_F^{\perp_M}(A)<\infty
 \quad\Longleftrightarrow\quad Y^\perp\vee F=\hat1 .
\]
Since $\rk(Y^\perp)=r-f$ and $\rk F=f$, modularity makes this
equivalent to $Y^\perp\wedge F=\hat0$.  Applying the lattice
anti-automorphism gives $Y\vee F^\perp=\hat1$, which is precisely the
basis condition for $M/F^\perp\oplus U_{0,F^\perp}$.  The finite values
are zero on both sides.  Thus the polar transform extends the given
flat involution.

\section*{Remaining open issues}

Two essential local-to-global inputs still need to be supplied before
the conditional theorem above becomes a complete proof.

First, Criterion \ref{crit:local} must be proved or replaced by a direct
verification of the full arbitrary incidence relations.  The invalid
shortcut is to apply Murota--Shioura to the union of two non-adjacent
basis layers; that union need not be an $M^\natural$-convex domain.
What is needed is either an induction from the nested exchange case to
the full Brandt--Eur--Zhang quotient exchange axiom, or a direct proof
of full incidence for the polar transform.

Second, the complete-flag local criterion used in Proposition
\ref{prop:set} must be stated with a precise non-uniform,
$\infty$-valued reference, or replaced by a direct proof that every
single polar layer $\nu^{\perp_M}$ is a valuated matroid.  The adjacent
opposition Lemma \ref{lem:elementary} proves the required cross
relations, but by itself it does not prove the Plucker relations inside
each transformed layer.

Modulo these two local criteria, the formula
$\nu\mapsto\nu^{\perp_M}$ is an order-reversing involution on
$\Dr(M)$ extending $F\mapsto F^\perp$.

\begin{thebibliography}{10}

\bibitem{BEZ}
M.~Brandt, C.~Eur and L.~Zhang,
\emph{Tropical flag varieties},
Advances in Mathematics \textbf{384} (2021), 107695.

\bibitem{DW}
A.~Dress and W.~Wenzel,
\emph{Valuated matroids},
Advances in Mathematics \textbf{93} (1992), 214--250.

\bibitem{DT}
A.~Dress and W.~Terhalle,
\emph{Well-layered maps---A class of greedily optimizable set functions},
Applied Mathematics Letters \textbf{8} (1995), 77--80.

\bibitem{JL}
M.~Jarra and O.~Lorscheid,
\emph{Flag matroids with coefficients},
Advances in Mathematics \textbf{436} (2024), 109396.

\bibitem{Murota}
K.~Murota,
\emph{Matroid valuation on independent sets},
Journal of Combinatorial Theory, Series B \textbf{69} (1997), 59--78.

\bibitem{MSsimple}
K.~Murota and A.~Shioura,
\emph{Simpler exchange axioms for $M$-concave functions on generalized
polymatroids},
Japan Journal of Industrial and Applied Mathematics \textbf{35} (2018),
235--259.

\bibitem{Speyer}
D.~Speyer,
\emph{Tropical linear spaces},
SIAM Journal on Discrete Mathematics \textbf{22} (2008), 1527--1558.

\end{thebibliography}

\end{document}
